\chapter{Weil Pairings and Automorphic Coefficient Multipliers}
\label{v4-arith:w5-a:chapter}

The Weil functional acts on compact smooth functions of the logarithmic
variable. Its zero formula evaluates their entire Fourier transforms.
Keeping these two spaces distinct fixes the prime weights and the
archimedean sign. It also distinguishes two coefficient multipliers,
one of which can produce only finite Fourier polynomials.

\section{The compact test algebra and its Fourier image}
\label{v4-arith:w5-a:weil-explicit}
\label{v4-arith:w5-a:weil-explicit:PW}

\begin{definition}[test spaces and involutions]
\label{v4-arith:w5-a:def:PW}
Put $\mathscr A=C_c^\infty(\mathbb R)$, with its usual
inductive-limit topology, and put $\mathscr P=\mathcal F_+\mathscr A$,
where
\[
 h(z)=\mathcal F_+f(z)=\int_{\mathbb R}f(t)e^{izt}dt,
 \qquad
 f(t)=\frac1{2\pi}\int_{\mathbb R}h(u)e^{-iut}du.
\]
Give $\mathscr P$ the transported topology. Its elements are entire
functions, rapidly decreasing in each horizontal strip.
For $f,g\in\mathscr A$, define
\[
 (f*g)(t)=\int f(v)g(t-v)dv,\qquad
 f^*(t)=\overline{f(-t)},\qquad h^\sharp(z)=\overline{h(\bar z)}.
\]
Thus $\mathcal F_+(f*g)=(\mathcal F_+f)(\mathcal F_+g)$ and
$\mathcal F_+(f^*)=(\mathcal F_+f)^\sharp$.
\end{definition}

The identities follow from Fubini and a change of variables. A second
real Fourier transform of $h\in\mathscr P$ is a reflected compact
smooth function. It is generally not an entire function available for
evaluation at complex zeros.

\begin{theorem}[the signed archimedean contribution]
\label{v4-arith:w5-a:thm:weil-explicit}
Set $\xi(s)=\tfrac12s(s-1)\pi^{-s/2}\Gamma(s/2)\zeta(s)$ and
$\psi_\Gamma=\Gamma'/\Gamma$.
For $f\in\mathscr A$, write $h=\mathcal F_+f$ and
$z_\rho=-i(\rho-1/2)$, where $\rho$ ranges over the zeros of $\xi$
with multiplicity.
Define
\begin{align}
 P(f)&=h(i/2)+h(-i/2),\notag\\
 W_\infty(f)&=\frac1{2\pi}\int_{\mathbb R}h(u)
       \bigl(\log\pi-\Re\psi_\Gamma(1/4+iu/2)\bigr)du,
 \label{v4-arith:w5-a:eq:weil-archimedean}\\
 W_{\mathrm{fin}}(f)&=\sum_{m\ge2}\frac{\Lambda(m)}{\sqrt m}
                   \bigl(f(\log m)+f(-\log m)\bigr).
 \label{v4-arith:w5-a:eq:weil-finite}
\end{align}
Then the continuous distribution $W:\mathscr A\to\mathbb C$ is
\begin{equation}
 W(f)=\sum_\rho h(z_\rho)=P(f)-W_\infty(f)-W_{\mathrm{fin}}(f).
 \label{v4-arith:w5-a:eq:weil-explicit}
\end{equation}
The zero sum converges absolutely and the prime sum is finite.
For even $f$, the prime sum is $2\mu_{1/2}(f)$.
\end{theorem}
\begin{proof}
Theorem~\ref{v4-arith:pm:weil} gives the formula with
$\mathcal A_\infty(h)=(2\pi)^{-1}\int h(\Re\psi_\Gamma-\log\pi)$.
The definition here is $W_\infty(f)=-\mathcal A_\infty(h)$.
This accounts for both signs in \eqref{v4-arith:w5-a:eq:weil-explicit}.
The compact-test estimates in that theorem prove the convergence and
continuity assertions. No tempered extension of $\mu_{1/2}$ is used.
\end{proof}

\begin{proposition}[a compact test determining the sign]
\label{v4-arith:wc:sign-test}
Let $0<a<b<\log2$. Suppose $f\ne0$ is even, nonnegative, smooth,
and supported in $(-b,-a)\cup(a,b)$. Then
\[
 W_{\mathrm{fin}}(f)=0,\qquad
 W_\infty(f)=2\int_a^b\frac{e^{-t/2}}{1-e^{-2t}}f(t)dt>0.
\]
\end{proposition}
\begin{proof}
There are no prime-power logarithms in the support. Since $f(0)=0$,
the constant terms of the digamma integral vanish after Fourier
inversion. The substitution $v=2t$ in
\[
 \psi_\Gamma(z)+\gamma
 =\int_0^\infty\frac{e^{-v}-e^{-zv}}{1-e^{-v}}dv
 \quad(\Re z>0)
\]
gives the displayed positive integral. The identity is
\href{https://dlmf.nist.gov/5.9.E16}{DLMF 5.9.16}.
The denominator is positive and $f$ is positive on an interval.
\end{proof}

\section{The Hermitian form and the two polar moments}
\label{v4-arith:w5-a:irreducible-obstruction}

Define $q_W(f,g)=W(f*g^*)$, with the first variable linear.
The zero multiset is invariant under $z\mapsto\bar z$, by the
functional equation and conjugation of $\xi$. Absolute convergence
therefore gives $W(f^*)=\overline{W(f)}$. Consequently $q_W$ is Hermitian.

\begin{corollary}[Weil's positivity criterion]
\label{v4-arith:w5-a:cor:weil-positivity-RH}
The Riemann Hypothesis is equivalent to
\begin{equation}
 q_W(f,f)\ge0\qquad(f\in\mathscr A).
 \label{v4-arith:w5-a:eq:weil-positivity}
\end{equation}
\end{corollary}
\begin{proof}
This is Weil's criterion on compact smooth additive tests, with the
Fourier convention above. It is stated in
\href{https://arxiv.org/html/2301.00421v3#S1}{Suzuki,
\emph{On the Hilbert space derived from the Weil distribution},
Section~1, equations~(1.1)--(1.2)}.
In that convention the zeros of $\xi(1/2-iz)$ are evaluated with
the opposite sign. Their reflected multiset is exactly $\{z_\rho\}$
here. Under RH the assertion follows directly from
$q_W(f,f)=\sum_\rho|\mathcal F_+f(\gamma_\rho)|^2$.
The converse is the cited positivity criterion, not an assertion of
positivity for this arithmetic form.
\end{proof}

\begin{proposition}[the polar Hermitian form]
\label{v4-arith:w5-a:rem:H3-removable}
Put $a(f)=\int f(t)e^{-t/2}dt$ and $b(f)=\int f(t)e^{t/2}dt$.
Then
\[
 P(f*g^*)=a(f)\overline{b(g)}+b(f)\overline{a(g)}.
\]
The moment map $(a,b):\mathscr A\to\mathbb C^2$ is continuous and
surjective. The polar form has rank two and both signs. It vanishes
on the common kernel of the two moments.
\end{proposition}
\begin{proof}
Evaluate $h_fh_g^\sharp$ at $i/2$ and $-i/2$.
If a linear combination of $a$ and $b$ vanished on every test, the
smooth function $c_1e^{-t/2}+c_2e^{t/2}$ would vanish identically.
Its value and derivative at zero force $c_1=c_2=0$. The moment map
therefore has two-dimensional range. On that range the form has
matrix $\left(\begin{smallmatrix}0&1\\1&0\end{smallmatrix}\right)$,
with eigenvalues $1$ and $-1$.
\end{proof}

The use of the entire function $\xi$ does not remove these polar
evaluations from its explicit formula. Removing a rank-two term
changes the form. It does not preserve positivity automatically.

\begin{conjecture}[archimedean Tate positivity]
\label{v4-arith:w5-a:conj:archimedean-tate-positivity}
For every $f\in\mathscr A$,
\begin{equation}
 W_\infty(f*f^*)\le P(f*f^*)-W_{\mathrm{fin}}(f*f^*).
 \label{v4-arith:w5-a:eq:archimedean-tate-positivity}
\end{equation}
\end{conjecture}

\begin{theorem}[the full positivity inequality]
\label{v4-arith:w5-a:thm:ATP-equals-weil-positivity}
\label{v4-arith:w5-a:thm:HPLi-restricted}
For all $f\in\mathscr A$,
\begin{equation}
 P(f*f^*)-W_{\mathrm{fin}}(f*f^*)=q_W(f,f)+W_\infty(f*f^*).
 \label{v4-arith:w5-a:eq:restricted-HPLi}
\end{equation}
Thus archimedean Tate positivity is equivalent to Weil positivity
and to RH. The equality applies to any specified restricted class
as well, but asserts no positivity on that class.
\end{theorem}
\begin{proof}
Apply \eqref{v4-arith:w5-a:eq:weil-explicit} to $f*f^*\in\mathscr A$
and rearrange. Then apply Weil's criterion.
\end{proof}

\section{Coefficient multipliers on the two carriers}
\label{v4-arith:w5-a:mellin-lift}

Fix a holomorphic cusp form $F(z)=\sum_{n\ge1}a_F(n)q^n$ of weight
$k>0$ for $\mathrm{SL}_2(\mathbb Z)$, where $k$ is even and $q=e^{2\pi iz}$.
Its coefficients have polynomial growth. Indeed the invariant
function $y^{k/2}|F(z)|$ is bounded on a fundamental domain and
hence on the upper half-plane. The Fourier coefficient integral
at $y=1/n$ gives $|a_F(n)|\le C n^{k/2}$.

\begin{definition}[the two coefficient maps]
\label{v4-arith:w5-a:def:upsilon}
For $f\in\mathscr A$ and $h=\mathcal F_+f\in\mathscr P$, set
\begin{equation}
 \begin{aligned}
 U_F(f)(z)&=\sum_{n\ge1}f(\log n)a_F(n)n^{(k-1)/2}q^n,\\
 \Upsilon_F(h)(z)&=\sum_{n\ge1}h(\log n)a_F(n)n^{(k-1)/2}q^n.
 \end{aligned}
 \label{v4-arith:w5-a:eq:upsilon}
\end{equation}
The first sum is finite. The second converges locally uniformly
on the upper half-plane, by polynomial coefficient growth and
the factor $e^{-2\pi n\Im z}$.
Neither definition asserts modularity.
\end{definition}

\begin{theorem}[the compact multiplier obstruction]
\label{v4-arith:w5-a:lem:upsilon-cuspidal}
The function $U_F(f)$ is a cusp form of weight $k$ for
$\mathrm{SL}_2(\mathbb Z)$ only if it vanishes identically.
Thus the compact multiplier is modular exactly on
$\{f:f(\log n)a_F(n)=0\text{ for all }n\}$.
The entire multiplier $\Upsilon_F(h)$ is a distinct construction.
Its modularity requires the transformation law of its resulting
holomorphic series.
\end{theorem}
\begin{proof}
Write $G(z)=\sum_{n=1}^N b_n e^{2\pi inz}$ for a nonzero finite
Fourier polynomial. Then $G(iy)$ is analytic as a function of $y$
near zero. If $G$ were modular of positive weight $k$,
\[
 G(i/y)=(iy)^kG(iy).
\]
The finite sum on the left is $O(e^{-2\pi/y})$ as $y\downarrow0$.
Hence $G(iy)=O(y^{-k}e^{-2\pi/y})$. Every Taylor coefficient at
zero must vanish, contradicting the nonzero analytic function.
This proves the first assertion. The coefficient condition is
exactly vanishing of the finite Fourier polynomial.
\end{proof}

\begin{remark}[Mellin domains]
\label{v4-arith:w5-a:rem:H2-obstruction}
\label{v4-arith:w5-a:rem:domain-PW-hol}
For $f\in\mathscr A$, the multiplicative Mellin integral
\[
 \int_0^\infty f(\log x)x^{s-1}dx=\int_{\mathbb R}f(t)e^{st}dt
\]
is entire and rapidly decreasing on every closed vertical strip.
This follows by integration by parts on the fixed compact support.
Strip holomorphy imposes no additional condition on this carrier.
For the entire multiplier $h(\log n)$, the analogous Mellin integral
has a different domain and must be specified separately.
\end{remark}

\section{The exact Rankin--Selberg coefficient calculation}
\label{v4-arith:w5-a:petersson}
\label{v4-arith:w5-a:mellin-lift:RS}

\begin{theorem}[unfolding for an actual cusp form]
\label{v4-arith:w5-a:thm:RS-integral}
Let $G(z)=\sum_{n\ge1}b_n e^{2\pi inz}$ be a holomorphic cusp form
of weight $k>0$ for $\mathrm{SL}_2(\mathbb Z)$.
Put $\Gamma=\mathrm{PSL}_2(\mathbb Z)$ and $d\mu=dx\,dy/y^2$.
Fix its standard fundamental domain $\mathcal F\subset\mathbb H$.
Use the Eisenstein series
$E(z,s)=\sum_{\Gamma_\infty\backslash\mathrm{PSL}_2(\mathbb Z)}
\Im(\gamma z)^s$, where $\Gamma_\infty$ is the translation subgroup.
For real $s>1$,
\begin{equation}
 \int_{\mathcal F}|G(z)|^2y^kE(z,s)d\mu
 =\frac{\Gamma(s+k-1)}{(4\pi)^{s+k-1}}
       \sum_{n\ge1}\frac{|b_n|^2}{n^{s+k-1}}.
 \label{v4-arith:w5-a:eq:KF-kernel}
\end{equation}
\end{theorem}
\begin{proof}
The Eisenstein series converges for $s>1$. Its cusp growth is
polynomial, while $G$ decreases exponentially, so the integral is
finite. Nonnegative summands permit unfolding by Tonelli:
the integral becomes
$\int_0^\infty\int_0^1|G(x+iy)|^2y^{s+k-2}dx\,dy$.
Fourier orthogonality gives $\sum|b_n|^2e^{-4\pi ny}$.
A second application of Tonelli and the gamma integral prove the formula.
\end{proof}

\begin{proposition}[positive norm and native pairing maps]
\label{v4-arith:w5-a:thm:petersson-positive}
The Petersson norm of an actual cusp form is positive definite.
If a complex vector space $E$ has an injective linear map
$I:E\to L^2(\mathcal F,y^k d\mu)$ on the fixed fundamental domain, then
\begin{equation}
 \langle v,w\rangle_I=\int_{\mathcal F}(Iv)(z)\overline{(Iw)(z)}y^k d\mu
 \label{v4-arith:w5-a:eq:petersson}
\end{equation}
is a positive definite form. Identifying another specified form
with this one requires the corresponding pairing identity for $I$.
\end{proposition}
\begin{proof}
These assertions follow from the norm property of $L^2$ and
injectivity. The unfolding formula computes the coefficients
of $G$ itself. For $G=\Upsilon_F(h)$ it therefore uses
$|h(\log n)a_F(n)n^{(k-1)/2}|^2$, whenever this series is modular.
\end{proof}

\begin{theorem}[a negative finite-prime square]
\label{v4-arith:w5-a:thm:prime-side}
There exists $f\in\mathscr A$ with $W_{\mathrm{fin}}(f*f^*)<0$.
Thus no linear map into a positive Hilbert space can realize
$W_{\mathrm{fin}}(f*g^*)$ as its inner product on all of $\mathscr A$.
\end{theorem}
\begin{proof}
Put $L=\log2$. Choose a nonzero real smooth $b$, supported in
$(-\varepsilon,\varepsilon)$, where
$2\varepsilon<\min(L,\log3-L)$. Write $b_L(t)=b(t-L)$ and
$f=b-b_L$. If $c=b*b^*$, then
\[
 f*f^*=2c-c(\,\cdot-L)-c(\,\cdot+L).
\]
The support meets the nonzero prime logarithms only at $L$ and
$-L$. At each, its value is $-c(0)=-\|b\|_2^2$. Therefore
\[
 W_{\mathrm{fin}}(f*f^*)=-\sqrt2\log2\,\|b\|_2^2<0.
\]
A Hilbert norm cannot have that value.
\end{proof}

\begin{corollary}[restricted automorphic comparisons]
\label{v4-arith:w5-a:cor:HPLi-restricted}
A positive realization of the finite-prime form on a restricted
subspace requires a proved restriction on that subspace.
The existence of an automorphic coefficient multiplier or an
interpolating function alone does not give a realization on
the full compact test algebra.
\end{corollary}
\begin{proof}
The preceding test excludes the full-space realization.
The Rankin--Selberg calculation identifies a different coefficient
kernel and supplies no equality with the displayed negative value.
\end{proof}

\section{Archimedean and interpolation obstructions}
\label{v4-arith:w5-a:sec:obstruction}

\begin{definition}[the archimedean density]
\label{v4-arith:w5-a:def:tate-archimedean}
The sign convention in \eqref{v4-arith:w5-a:eq:weil-explicit} uses
\begin{equation}
 \mathcal T_\infty(u)=
 \frac{\log\pi-\Re\psi_\Gamma(1/4+iu/2)}{2\pi}.
 \label{v4-arith:w5-a:eq:tate-archimedean-distribution}
\end{equation}
\end{definition}

\begin{proposition}[the gamma density is not the Weil density]
\label{v4-arith:w5-a:prop:archimedean-not-absorbed}
The functions $|\Gamma(1/4+iu/2)|^2$ and $\mathcal T_\infty(u)$
are not constant multiples. In fact $\mathcal T_\infty$ has both signs.
For a local continuous argument of the nonvanishing gamma function,
\[
 \frac{d}{du}\arg\Gamma(1/4+iu/2)=\frac12\Re\psi_\Gamma(1/4+iu/2),
 \qquad
 \frac{d}{du}\log|\Gamma(1/4+iu/2)|^2=-\Im\psi_\Gamma(1/4+iu/2).
\]
\end{proposition}
\begin{proof}
Differentiate $\log\Gamma(1/4+iu/2)$ to obtain $(i/2)\psi_\Gamma$.
Taking real and imaginary parts proves the identities.
The digamma series gives $\psi_\Gamma(1/4)<-\gamma<0$.
For $y\ge1$, its real part is
\[
 -\gamma+\sum_{n\ge0}
 \left(\frac1{n+1}-\frac{n+1/4}{(n+1/4)^2+y^2}\right).
\]
The first $\lfloor y\rfloor$ terms have sum at least $\log y-C$:
the harmonic part has this growth, and the subtracted part is bounded
by $y^{-2}\sum_{n\le y}(n+1/4)\le C$.
In the remaining terms the possible negative part is at most
$C\sum_{n>y}n^{-2}$. Thus the real part tends to positive infinity.
The density $\mathcal T_\infty$ is positive at zero and negative
for sufficiently large $|u|$. The gamma absolute square is positive.
\end{proof}

\begin{proposition}[rational Li tests are not compact Fourier tests]
\label{v4-arith:w5-a:prop:li-via-weil}
For any integer $n\ge1$, there is no $f\in\mathscr A$ such that
\[
 (\mathcal F_+f)(-i(\rho-1/2))=1-(1-1/\rho)^n
\]
for every nontrivial zero $\rho$.
\end{proposition}
\begin{proof}
The right side is $n/\rho+O_n(|\rho|^{-2})$ as $|\rho|\to\infty$.
The left side decreases faster than every inverse power of
$|\Im\rho|$, uniformly in the zero strip. There are infinitely many
zeros and only finitely many in any bounded set. Since their real
parts lie in $[0,1]$, these two estimates contradict one another.
\end{proof}

\begin{proposition}[continuous-spectrum pairing data]
\label{v4-arith:w5-a:prop:eisenstein-residual}
Let $\int^\oplus E_u\,d\mu(u)$ be a specified Hilbert direct integral.
An identity between its norm and an arithmetic Hermitian form requires
a map into that direct integral and equality of the two pairings.
A scalar gamma factor determines neither that map nor its measure.
\end{proposition}
\begin{proof}
The direct-integral norm is $\int\|v(u)\|_{E_u}^2d\mu(u)$.
Replacing a nonzero realization map $I$ by $2I$ multiplies its
pullback pairing by four. A single scalar factor does not select
the map or its measure. The sign-changing density above also
cannot equal an unspecified positive norm by scalar rescaling.
\end{proof}

\section{The arithmetic pairing requirement}
\label{v4-arith:w5-a:form}

\begin{theorem}[the remaining positive comparison]
\label{v4-arith:w5-a:thm:hpli-form}
A Hilbert realization of the full Weil form consists of a positive
Hilbert space $K$, a linear map $J:\mathscr A\to K$, and the identity
\[
 \langle Jf,Jg\rangle_K=W(f*g^*)\qquad(f,g\in\mathscr A).
\]
Such a realization exists if and only if $q_W$ is positive semidefinite.
When it exists, its minimal realization is the completion of
$\mathscr A/\{f:q_W(f,f)=0\}$. Its existence is RH-equivalent.
Neither coefficient map in \eqref{v4-arith:w5-a:eq:upsilon}
provides the pairing identity.
\end{theorem}
\begin{proof}
A Hilbert realization implies nonnegativity. Conversely, positivity
and polarization give Cauchy--Schwarz for $q_W$. Its nullspace is a
linear subspace orthogonal to every vector. Quotienting and completing
therefore gives the asserted Hilbert space and canonical map.
Weil's criterion supplies the equivalence to RH. The coefficient
maps construct holomorphic functions with specified coefficients.
Their definitions contain no equality with the full Weil form, and
the finite-prime counterexample excludes the proposed full-space
positive realization of that separate term.
\end{proof}

The arithmetic problem asks for this pairing through a specified
native carrier, its operators, and its realization maps. The compact
formula, the exact coefficient calculation, and the two obstruction
theorems retain those separate requirements.
